\section{Практическая часть}

\subsection{Определение радиуса кривизны поверхности линзы}
\begin{table}[H]
	\centering
	\caption{Таблица измерений границ колец Ньютона с микрометра}
	\begin{tabular}{cccc}
		\toprule
		№ &   $l_1$, мм &  $l_2$, мм &  $l_3$, мм \\
		\midrule
		10 &   31,20 &  31,20 &  31,28 \\
		9 &   31,43 &  31,43 &  31,43 \\
		8 &   31,59 &  31,57 &  31,58 \\
		7 &   31,71 &  31,72 &  31,71 \\
		6 &   31,86 &  31,90 &  31,89 \\
		5 &   32,02 &  32,04 &  32,02 \\
		4 &   32,22 &  32,26 &  32,21 \\
		3 &   32,45 &  32,47 &  32,74 \\
		2 &   32,73 &  32,74 &  32,77 \\
		1 &   33,09 &  33,08 &  33,07 \\
		-1 &   34,62 &  34,63 &  34,61 \\
		-2 &   34,99 &  35,02 &  35,02 \\
		-3 &   35,26 &  35,29 &  35,27 \\
		-4 &   35,50 &  35,56 &  35,50 \\
		-5 &   35,69 &  35,75 &  35,69 \\
		-6 &   35,87 &  35,91 &  35,89 \\
		-7 &   36,03 &  36,06 &  36,04 \\
		-8 &   36,18 &  36,19 &  36,17 \\
		-9 &   36,34 &  36,36 &  36,33 \\
		-10 &   36,48 &  36,49 &  36,49 \\
		\bottomrule
	\end{tabular}
\end{table}

Для определения диаметров колец найдем разность значений номеров положительных и отрицательных (10 и -10, 9 и -9, ...), полученные значения занесем в \cref{tab:2}. Для построения графика квадрата диаметров колец от их номера составим таблицу с вычисленными значениями и погрешностями и их средних значений. 
\begin{table}[H]
	\centering
	\caption{Диаметры колец Ньютона }
	\label{tab:2}
	\begin{tabular}{cccc}
		\toprule
		$k$ &  $d_{k\, 1}$, мм & $d_{k\, 2}$, мм & $d_{k\, 3}$, мм \\
		\midrule
		10 &  5,28 & 5,29 & 5,21 \\
		9  &  4,91 & 4,93 & 4,90 \\
		8  &  4,59 & 4,62 & 4,59 \\
		7  &  4,32 & 4,34 & 4,33 \\
		6  &  4,01 & 4,01 & 4,00 \\
		5  &  3,67 & 3,71 & 3,67 \\
		4  &  3,28 & 3,30 & 3,29 \\
		3  &  2,81 & 2,82 & 2,53 \\
		2  &  2,26 & 2,28 & 2,25 \\
		1  &  1,53 & 1,55 & 1,54 \\
		\bottomrule
	\end{tabular}
\end{table}
\begin{table}[H]
	\centering
	\caption{Квадраты диаметров колец Ньютона и их погрешность}
	\begin{tabular}{ccccccccc}
		\toprule
		$k$ & \multicolumn{3}{c}{$d_k^2$, мм} & \multicolumn{3}{c}{$\Delta d_k^2 \cdot 10^{-5}$, мм} &  $\overline{d}_k^2$, мм & $\Delta \overline{d}_k^2 \cdot 10^{-5}$, мм\\
		\midrule
		& 1 & 2 & 3 & 1 & 2 & 3 & & \\
		\midrule
		10 & 27,88 & 27,98 & 27,14 & 5,28 & 5,29 & 5,21 & 27,67& 3,04\\
		9  & 24,11 & 24,30 & 24,01 & 4,91 & 4,93 & 4,90 & 24,14& 2,84\\
		8  & 21,07 & 21,34 & 21,07 & 4,59 & 4,62 & 4,59 & 21,16& 2,66\\
		7  & 18,66 & 18,84 & 18,75 & 4,32 & 4,34 & 4,33 & 18,75& 2,50\\
		6  & 16,08 & 16,08 & 16,00 & 4,01 & 4,01 & 4,00 & 16,05& 2,31\\
		5  & 13,47 & 13,76 & 13,47 & 3,67 & 3,71 & 3,67 & 13,57& 2,13\\
		4  & 10,76 & 10,89 & 10,82 & 3,28 & 3,30 & 3,29 & 10,82& 1,90\\
		3  & 7,90  & 7,95  & 7,84  & 2,81 & 2,82 & 2,80 & 7,90& 1,62\\
		2  & 5,11  & 5,20  & 5,06  & 2,26 & 2,28 & 2,25 & 5,12& 1,31\\
		1  & 2,34  & 2,40  & 2,37  & 1,53 & 1,55 & 1,54 & 2,37& 0,90\\
		\bottomrule& 
	\end{tabular}
\end{table}

Найдем радиусы линзы и погрешность по формулам \cref{eq:th:1,eq:th:2} соответственно при параметрах:
\[
n = 1 \qquad 
\lambda = 0,547\ \text{мкм} \qquad
\Delta d_k = 5\ \text{мкм} \qquad
\beta = 22,0^\circ \qquad
\Delta \beta = 0,5^\circ
\]
\begin{table}[H]
	\centering
	\caption{Результаты вычисления радиуса линзы и ее погрешности}
	\begin{tabular}{cccccccc}
		\toprule
		$k$ & $R_1$, мм & $R_2$, мм & $R_3$, мм & $\Delta R_1$, мм & $\Delta R_2$, мм & $\Delta R_3$, мм \\
		\midrule
		10 & 921,24 & 945,48 & 933,32 & 192,12 & 197,10 & 194,60 \\
		9  & 1005,03 & 1022,89 & 996,15 & 207,48 & 211,12 & 205,66 \\
		8  & 1035,81 & 1043,20 & 839,67 & 212,93 & 214,44 & 172,94 \\
		7  & 1058,47 & 1071,42 & 1064,93 & 217,05 & 219,69 & 218,37 \\
		6  & 1060,11 & 1083,35 & 1060,11 & 217,05 & 221,77 & 217,05 \\
		5  & 1054,70 & 1054,70 & 1049,44 & 215,69 & 215,69 & 214,623 \\
		4  & 1049,20 & 1058,94 & 1054,07 & 214,38 & 216,36 & 215,37 \\
		3  & 1036,40 & 1049,99 & 1036,40 & 211,62 & 214,38 & 211,62 \\
		2  & 1054,17 & 1062,78 & 1049,88 & 215,10 & 216,85 & 214,23 \\
		1  & 1097,13 & 1101,29 & 1068,23 & 223,71 & 224,56 & 217,85 \\
		\bottomrule
	\end{tabular}
\end{table}
